%/////////////////////////////////////////////////////////////////////
%	 LRodrigo - Cores e Comando
%	Ultima atualização - 27/02/2018
%/////////////////////////////////////////////////////////////////////




%%%%%%%%%%%%%%%%%%%%%%%%%%%%%%%%%%%%%%%%%
\newcommand\ReduzirFonte{\fontsize{6}{7.2}\selectfont}


%/////////////////////////////////////////////////////////////////////
%      Tipos de Colunas
%/////////////////////////////////////////////////////////////////////
\newcolumntype{G}{>{\columncolor{gray!25}}c}
\newcolumntype{g}{>{\columncolor{gray!25}}c}
\newcolumntype{B}{>{\columncolor{bluegray}}c}

\definecolor{airforceblue}{rgb}{0.36, 0.54, 0.66}
%\newcolumntype{b}{>{\columncolor{airforceblue}}c}

\definecolor{beige}{rgb}{0.96, 0.96, 0.86}
%\newcolumntype{b}{>{\columncolor{beige}}c}


%/////////////////////////////////////////////////////////////////////
% Cor da celula de destque
%/////////////////////////////////////////////////////////////////////
\definecolor{antiquebrass}{rgb}{0.8, 0.58, 0.46}
\newcolumntype{A}{>{\columncolor{antiquebrass}}c}
\newcolumntype{a}{>{\columncolor{antiquebrass}}l}


%/////////////////////////////////////////////////////////////////////
% Azul Centralizado
%/////////////////////////////////////////////////////////////////////
\definecolor{aliceblue}{rgb}{0.94, 0.97, 1.0}
\newcolumntype{b}{>{\columncolor{aliceblue}}c}

\definecolor{beaublue}{rgb}{0.74, 0.83, 0.9}
\newcolumntype{d}{>{\columncolor{beaublue}}c}

\definecolor{babyblueeyes}{rgb}{0.63, 0.79, 0.95}
\newcolumntype{e}{>{\columncolor{babyblueeyes}}c}

\definecolor{babyblue}{rgb}{0.54, 0.81, 0.94}
\newcolumntype{f}{>{\columncolor{babyblue}}c}




%/////////////////////////////////////////////////////////////////////
% Background personalizado
%		As imagems que podem ser utilizadas como background devem
%	ser armazenadas no diretorio "Feathergraphics"
%/////////////////////////////////////////////////////////////////////
\newcommand{\backGround}[1]{
	\setbeamertemplate{background}
	{
		\includegraphics[width=\paperwidth,height=\paperheight]{Feathergraphics/{#1}}
		\tikz[overlay] \fill[fill opacity=0.75,fill=white] (0,0) rectangle (-\paperwidth,\paperheight);
	}
}


%/////////////////////////////////////////////////////////////////////
% Alterando o simbolo utilizado no timemize
%		\bullet
%		\ast
%		\bigcirc
%		\circ
% Lista de Simbolos:
%	http://tug.ctan.org/info/symbols/comprehensive/symbols-a4.pdf
%   https://docs.latexbase.com/symbols/
%/////////////////////////////////////////////////////////////////////
%\mode<presentation>
%{
%	%\usetheme{Madrid}
%	\usefonttheme{professionalfonts} 
%	%\setbeamertemplate{itemize item}{\color{airforceblue}$\blacksquare$}
%	\setbeamertemplate{itemize item}{\color{airforceblue}$\blacktriangleright$}
%	
%	\setbeamertemplate{itemize subitem}{\color{airforceblue}$\checkmark$}
%
%}

%/////////////////////////////////////////////////////////////////////
% Outra forma de trocar os simbolos
%/////////////////////////////////////////////////////////////////////
%\setbeamercolor{itemize item}{fg=red}
%\setbeamercolor{itemize subitem}{fg=blue}
%\setbeamercolor{itemize subsubitem}{fg=cyan}
%
%\setbeamertemplate{itemize item}[square]
%\setbeamertemplate{itemize subitem}[circle]
%\setbeamertemplate{itemize subsubitem}[triangle]
